\subsection{DET-15: actual convolution projections on one Bessel \texorpdfstring{$L^2$}{L2} space}
There is a stronger common-space construction in the source than the
finite target of \fid{DET-14}. Its matrix-coefficient convolution is
an orthogonal rank-one projection on a Bessel \(L^2\) model.
This gives positive finite operators of arbitrarily large rank on
fixed side data. It also permits an exact realization of the
zero-based operator \(B\) in that model. The last realization uses
the zero-dependent weights explicitly; it does not identify an
independent arithmetic choice of weights or the historical
determinant family.

\paragraph{Exact source hypotheses and normalization.}
Put \(G=\mathrm{UH}(H_{k,r})\), \(S=S_{k,r,z}\),
\(\psi=\mathrm{psiGen}(k,r,z,c)\), and let \(W_p\) be the finite-dimensional
Hilbert space in the side representation \(\pi\).
The definitions in
\path{Development/AbelianVarietyTheory/Definitions/SemilocalFrames/MatrixOfBesselTwoSidedAction.lean}
are as follows. Lines 47--48 define
\[
 D=\mathrm{DD}(\sigma,w_0)
 =\operatorname{span}\{\sigma(g)w_0:g\in G\}.
\]
Unless a different directory is displayed, the abbreviated Bessel
proof filenames below are in
\path{Development/AbelianVarietyTheory/Proofs/FlagMeasures/};
the abbreviated \path{BesselPeriodGGP.lean} remains the file in
\path{Formalism/PeriodIntegrals/}.
The condition \texttt{Std}, lines 53--55, includes \(k\ge1\), a
strongly continuous unitary irreducible representation \(\sigma\)
on a complete Hilbert space \(H_\sigma\), \(\|w_0\|=1\), a nonzero
left-invariant sigma-finite measure \(\mu_G\) supported on \(G\),
the Schur relation \texttt{SchurH}, a specified lower bound on
the decay exponent \(a\), and \texttt{Decay} on \(D\).
In particular, \texttt{SchurH}, lines 42--45, says \(d>0\) and
\[
 \int_G
 \langle v,\sigma(g)u\rangle
 \overline{\langle v',\sigma(g)u'\rangle}\,d\mu_G(g)
 =\frac{\langle u',u\rangle\langle v,v'\rangle}{d},
\]
with the indicated products integrable.
The condition \texttt{Side}, lines 60--62, includes the exact
support and \(T(z)>0\) conditions, a unitary \(\pi\) on the
specified compact subgroup, a nonzero left-invariant sigma-finite
\(\mu_S\) on \(S\), and finite-dimensional \(W_p\).
We retain all these hypotheses, as well as a measure \(\nu\)
satisfying \texttt{IsQuotMeas}\((\mu_G,\mu_S,\nu)\).

The condition \texttt{IsEll}, lines 119--121, supplies a Bessel map
\(l:H_\sigma\to W_p\), a vector \(v_0\in D\) with \(l(v_0)\ne0\),
and
\[
 \operatorname{be}(u,v)=l(u)l(v)^*,\qquad u,v\in D.
\]
This normalization is a retained hypothesis. An actual existence
export is
\path{Development/AbelianVarietyTheory/Proofs/FlagMeasures/ExistsBesselMapRankOneFactorisationTranslateSpanStatement.lean},
lines 7--8: it supplies such an \(l\) from \texttt{Std},
\texttt{Side}, smoothness of \(w_0\), the stated multiplicity-one
condition, the quotient measure, and the stated nonvanishing
condition. We do not omit those assumptions.

The actual definitions in
\path{Formalism/PeriodIntegrals/BesselPeriodGGP.lean}, lines 79--92,
are
\[
 (\operatorname{rconv}(e)W)(x)
 =\int_G e(g)W(xg)\,d\mu_G(g),\qquad
 e_u(g)=\mathrm{fdegH}(\mu_G,\sigma,u)
             \overline{\langle u,\sigma(g)u\rangle}.
\]
The Schur relation gives
\(\mathrm{fdegH}(\mu_G,\sigma,u)=d\) for every nonzero \(u\):
its denominator is \(\|u\|^4/d\).
This is also the precise export
\path{Development/AbelianVarietyTheory/Proofs/FlagMeasures/BesselFormalDegreeEqOfSchurOrthogonalityStatement.lean},
lines 7--8, with its complete elementary proof in the corresponding
\path{BesselFormalDegreeEqOfSchurOrthogonality.lean}, lines 7--29.
Thus for every unit \(u\) used below,
\[
 e_u(g)=d\,\overline{\langle u,\sigma(g)u\rangle};
\]
there is no additional normalization by \(\|u\|\).

\paragraph{The common Hilbert space and retained measurable domain.}
Let \(\mathcal D\) be the vector space of functions
\(W:\mathrm{Mat}\to W_p\) which are almost everywhere strongly
measurable for both \(\mu_G\) and \(\nu\), satisfy the source's
\(S,\psi,\pi\)-equivariance, and have finite \texttt{l2SG} norm
for \(\nu\). The equivariance and norm definitions are
\path{BesselPeriodGGP.lean}, lines 68--77.
Let
\[
 \mathcal B=\overline{\{[W]_\nu:W\in\mathcal D\}}
       \ \subset L^2(\nu;W_p).
\]
This closed Hilbert subspace depends only on the fixed side and
measure data, not on the representation vector to be selected.
We use closure here; no unproved density in all of \(L^2(\nu;W_p)\)
is asserted.

The source's quotient-measure theorem
\path{Development/AbelianVarietyTheory/Proofs/FlagMeasures/WeilFormulaBesselUnitaryIntegrableStatement.lean},
lines 7--8, gives sigma-finiteness of \(\nu\), its support on \(G\),
and the integrable Weil formula, under its closed-subgroup,
Haar and unimodularity hypotheses. Those group hypotheses in
this setting are the ones used explicitly in
\path{BesselFunctionalStandardSide.lean}, lines 59--73.
Support on \(G\) is also proved directly in
\path{BesselModelProjectorOnLine.lean}, lines 40--69.

\paragraph{DET-15(a): the actual convolution is a positive rank-one projection.}
Under the retained source hypotheses, put
\[
 f_u(g)=l(\sigma(g)u),\qquad u\in D,\qquad
 \eta_u=\sqrt d\,[f_u]_\nu.
\]
These classes belong to \(\mathcal B\), and
\[
 \langle\eta_v,\eta_u\rangle_{\mathcal B}
 =\langle v,u\rangle_{H_\sigma}.
\]
For \(u=w_0\), the actual convolution on \(\mathcal D\) extends
uniquely to the orthogonal rank-one projection
\[
 P_{w_0}W=\langle\eta_{w_0},W\rangle\eta_{w_0}
 \quad\text{on }\mathcal B.
\]
It has norm and trace \(1\), is positive and trace class, and its
ordinary Fredholm determinant is \(1-w\).
\begin{proof}
We first establish the measurable representatives and the full
inner-product normalization. Choose \(v_0\in D\) with
\(b_0=l(v_0)\ne0\), and put \(N=\|b_0\|^2>0\).
Let
\[
 A_u(g)=\operatorname{be}(\sigma(g)u,v_0)b_0.
\]
The space \(D\) is invariant under \(G\), since it is the span
of an orbit. Thus the \texttt{IsEll} factorization gives, on \(G\),
\[
 A_u(g)=N f_u(g).
\]
The source theorem
\path{Development/AbelianVarietyTheory/Proofs/FlagMeasures/NormIdentityBesselCoefficientQuotientStatement.lean},
lines 7--8, makes \(A_u\) equivariant and in \(L^2(\nu;W_p)\);
it also supplies the mixed inner-product identity below.
The Bessel-form continuity export
\path{BesselFormConvergesEquivariantHermitianStatement.lean},
lines 7--8, applied with the invariant domain \(D\) and its
retained decay bound, makes \(A_u\) continuous on \(G\).
This application and the ambient-measurability step are displayed
explicitly in
\path{LintegralBesselCoefficientCompactLe.lean}, lines 29--53.
Since \(G\) is closed in the finite-dimensional matrix space,
the zero extension of \(A_u/N\) off \(G\) is Borel measurable
and hence strongly measurable for both measures, since \(W_p\)
is finite dimensional. On \(G\) it is \(f_u\), and both measures
are supported on \(G\). This gives a representative in
\(\mathcal D\) for \([f_u]_\nu\); arbitrary values of \(l\)
off this set cause no ambiguity.

The same norm-identity export gives
\[
 \int\langle A_v(g),A_u(g)\rangle\,d\nu(g)
 =\frac{\langle v,u\rangle
       \langle b_0,\operatorname{be}(v_0,v_0)b_0\rangle}{d}
 =\frac{N^2\langle v,u\rangle}{d}.
\]
The left side is \(N^2\langle[f_v],[f_u]\rangle\).
Cancelling \(N^2>0\) proves the claimed isometry identity.
It also proves complex linearity of \(u\mapsto\eta_u\):
the squared norm of the image of any finite linear relation,
expanded using this identity, is the squared norm of that
relation in \(H_\sigma\), hence zero when the relation is zero.
No global boundedness of the point-evaluation map \(l\) is assumed.

The exact convolution export is
\path{Development/AbelianVarietyTheory/Proofs/FlagMeasures/BesselModelProjectorOnLineStatement.lean},
lines 7--8. Its complete proof is
\path{BesselModelProjectorOnLine.lean}, lines 71--137.
For every \(W\in\mathcal D\), it states, for every \(y\in G\),
\[
 \operatorname{rconv}(e_{w_0})W(y)
 =d\left(\int\langle f_{w_0}(g),W(g)\rangle\,d\nu(g)\right)
      f_{w_0}(y).
\]
Its integrable-kernel step is the source export
\path{BesselFunctionalStandardSideStatement.lean}, lines 7--8;
the full underlying proof is
\path{BesselFunctionalStandardSide.lean}, lines 57--82.
The scalar integral is absolutely convergent by Cauchy--Schwarz,
because both factors lie in \(L^2(\nu;W_p)\).
As \(\nu\) is supported on \(G\), passage to equivalence classes
identifies this convolution with
\(\langle\eta_{w_0},[W]\rangle\eta_{w_0}\).
In particular, it depends only on \([W]_\nu\).

The isometry identity and \(\|w_0\|=1\) give
\(\|\eta_{w_0}\|=1\).
For a unit vector \(\eta\), the map \(W\mapsto\langle\eta,W\rangle\eta\)
is bounded of norm \(1\), selfadjoint, idempotent and positive,
since its quadratic form is \(|\langle\eta,W\rangle|^2\).
Its range is the line \(\mathbb C\eta\), so its sole nonzero
eigenvalue is \(1\), proving the trace and determinant.
The source convolution agrees with this bounded operator on the
dense subspace defining \(\mathcal B\), which proves its unique
bounded extension.
\end{proof}

\paragraph{DET-15(b): the fixed model has infinitely many orthogonal source projections.}
The retained \texttt{Std} hypotheses imply that \(H_\sigma\)
is infinite dimensional. There exists an orthonormal sequence
\(u_1,u_2,\ldots\) in \(D\) such that, with the same side data,
measure, representation, decay exponent, formal degree and map \(l\),
each actual convolution \(\operatorname{rconv}(e_{u_j})\)
extends to \(P_j=\eta_{u_j}\eta_{u_j}^*\).
The vectors \(\eta_{u_j}\) are orthonormal in \(\mathcal B\).
\begin{proof}
First, \(\mu_G(G)=\infty\); this can be proved without an
unstated Radon or local-finiteness assumption. Since \(k\ge1\),
the middle block of \(H_{k,r}\) has the positive coordinate \(r\)
and negative coordinate \(r+k\). Let \(g(t)\) act as the
real Lorentz boost on these two coordinates and as the identity
on all others. Then \(g(t)\in G\). The nonzero row covector
\(\ell=e_r^T+e_{r+k}^T\) satisfies \(\ell g(t)=e^t\ell\).
Every \(x\in G\) is invertible, since \(H_{k,r}\) is invertible,
so \(Q(x)=\|\ell x\|>0\). The measurable sets
\[
 E_n=\{x\in G:e^n\le Q(x)<e^{n+1}\},\qquad n\in\mathbb Z,
\]
partition \(G\), and \(g(1)E_n=E_{n+1}\).
Left invariance gives the same measure to all these sets.
If \(\mu_G(G)\) were finite, their common measure would be zero,
and countable additivity would force \(\mu_G(G)=0\).
That contradicts the nonzero measure supported on \(G\).

If \(H_\sigma\) had finite dimension \(D_0\), then \(D_0\ge1\)
because \(\|w_0\|=1\). For an orthonormal basis \(e_1,\ldots,e_{D_0}\),
unitarity gives pointwise on \(G\)
\[
 \sum_{i,j=1}^{D_0}
 |\langle e_i,\sigma(g)e_j\rangle|^2=D_0.
\]
The Schur relation integrates every summand to \(1/d\).
Its finite sum would therefore give
\(D_0\mu_G(G)=D_0^2/d<\infty\), a contradiction.
Thus \(H_\sigma\) is infinite dimensional.

The closure of \(D\) is nonzero and invariant under every bounded
\(\sigma(g)\). Irreducibility therefore makes it all of \(H_\sigma\).
Recursively choose a vector in \(D\) outside the span of the previously
chosen orthonormal vectors, subtract its finite orthogonal projection
onto that span, and normalize. Each result still lies in the
linear space \(D\), and the recursion never stops since \(D\)
is dense in the infinite-dimensional \(H_\sigma\).
This produces the asserted sequence. The isometry in part (a)
makes the corresponding \(\eta_{u_j}\) orthonormal.

It remains to check that each projection is the stated source
convolution, rather than only an abstract projection. For any unit
\(u\in D\), the orbit span \(D_u=\mathrm{DD}(\sigma,u)\) is
contained in \(D\) by invariance. All clauses of \texttt{Std}
remain unchanged except the unit vector and its orbit span;
the unit norm holds, and the decay bound restricts from \(D\)
to \(D_u\). Thus \texttt{Std}\((\sigma,u,a,\mu_G,d)\) holds.
The \texttt{IsEll} factorization restricts to \(D_u\).
Its nonvanishing clause also holds: part (a) gives
\(\int\|l(\sigma(g)u)\|^2\,d\nu(g)=1/d>0\).
As \(\nu\) is supported on \(G\), some \(g\in G\) has
\(l(\sigma(g)u)\ne0\), and \(\sigma(g)u\in D_u\).
The Bessel-map clause itself is unchanged.
All hypotheses of the source projector theorem therefore hold
with \(w_0\) replaced by \(u\), proving the claim.
\end{proof}

\paragraph{DET-15(c): explicit trace-class convolution sums.}
Let \((\beta_j)\) be any complex sequence with
\(\sum_j|\beta_j|<\infty\). On this one fixed Hilbert space,
\[
 T_N=\sum_{j\le N}\beta_jP_j,\qquad
 T_\beta=\sum_{j\ge1}\beta_jP_j
\]
are the finite and limiting operators. For each finite \(N\),
\(T_N\) is the bounded extension of the actual source convolution
with \(e_N=\sum_{j\le N}\beta_j e_{u_j}\) on \(\mathcal D\).
Moreover,
\[
 \|T_\beta-T_N\|_1=\sum_{j>N}|\beta_j|,\qquad
 \Tr T_\beta=\sum_j\beta_j,\qquad
 \det(I-wT_\beta)=\prod_j(1-w\beta_j).
\]
The operator is normal; it is positive exactly when every
\(\beta_j\) is real and nonnegative.
\begin{proof}
The orthonormal vectors \(\eta_{u_j}\) diagonalize these operators
on their closed span, and the orthogonal complement is killed.
Consequently the singular values are \(|\beta_j|\), including
multiplicity, and the stated trace-norm, trace and determinant
formulas follow from the absolutely convergent diagonal series
and product. Normality follows by the same diagonal computation.
Positivity is sufficient for nonnegative real \(\beta_j\), and
necessary by testing the quadratic form on each \(\eta_{u_j}\).
For finite \(N\), linearity of the finitely many integrable
convolution expressions on \(\mathcal D\), followed by part (a),
identifies \(T_N\) with the stated convolution.
\end{proof}
For nonnegative weights, the exact trace tail is \(\sum_{j>N}\beta_j\).
The unweighted partial sums instead have trace \(N\), so they
cannot converge in trace norm. Thus the common-space problem for
these actual projections is solved; the trace-class passage
requires a specified summable weighting.
For example, \(\beta_j=2^{-j}\) gives a positive trace-class
construction independent of all zeros, with trace \(1\) and
exact trace-norm tail \(2^{-N}\). This example does not claim
the Xi determinant.

\paragraph{An actual measurable integral kernel.}
Write \(\eta_j(x)=\sqrt d\,f_{u_j}(x)\) using the measurable
representatives constructed above. Then
\[
 K_\beta(x,y)=\sum_j\beta_j\eta_j(x)\eta_j(y)^*
\]
converges absolutely in operator norm for almost every \((x,y)\),
belongs to \(L^2(\nu\times\nu;\mathrm{HS}(W_p))\), and is an
integral kernel for \(T_\beta\) on \(\mathcal B\).
\begin{proof}
Put \(a(x)=\sum_j|\beta_j|\|\eta_j(x)\|^2\).
Tonelli and \(\|\eta_j\|_2=1\) give
\(\int a\,d\nu=\sum_j|\beta_j|<\infty\), so \(a(x)<\infty\)
almost everywhere. Cauchy--Schwarz for the series gives
\[
 \sum_j|\beta_j|\|\eta_j(x)\|\,\|\eta_j(y)\|
 \le\sqrt{a(x)a(y)}.
\]
The rank-one matrix norm and Hilbert--Schmidt norm are both
\(\|\eta_j(x)\|\,\|\eta_j(y)\|\).
This proves absolute convergence and measurability, and bounds
the squared Hilbert--Schmidt norm of the sum by \(a(x)a(y)\),
which is integrable on the product. Applying the same bound
to tails proves convergence in product \(L^2\).

For almost every fixed \(x\) and any \(W\in L^2(\nu;W_p)\),
\[
 \int\|K_\beta(x,y)W(y)\|\,d\nu(y)
 \le\sqrt{a(x)}\left(\sum_j|\beta_j|\right)^{1/2}\|W\|_2,
\]
again by Cauchy--Schwarz. Thus the kernel integral is defined.
For an \(L^2\) matrix kernel, pointwise Cauchy--Schwarz and
integration show that the operator norm is at most its
product \(L^2\) Hilbert--Schmidt norm. Therefore the finite
kernel sums converge as operators to this integral operator.
Their operators are precisely \(T_N\), so the limit is
\(T_\beta\), including on \(\mathcal B\).
\end{proof}

\paragraph{Exact embedding of the zero-based operator, with its dependence exposed.}
Let \(\gamma_j\) be the one-per-\(\pm\)-pair Xi-zero list, with
multiplicity, from \fid{B-01}, and take
\(\beta_j=\gamma_j^{-2}\).
No RH is needed for \(\sum_j|\beta_j|<\infty\).
The isometric embedding
\[
 V:\ell^2\longrightarrow\mathcal B,\qquad
 Ve_j=\eta_{u_j},
\]
then gives \(T_\beta=VBV^*\). On
\(V\ell^2\oplus(V\ell^2)^\perp\), this is \(B\oplus0\), and
\[
 \det_{\mathcal B}(I-wT_\beta)=G(w)
 =\frac{\Xi(\sqrt w)}{\Xi(0)}.
\]
These claims follow directly from part (c) and \fid{B-01};
the determinant receives factor \(1\) from the killed complement.
The operator is positive if and only if RH, by the same
eigenvalue criterion as \(B\). The cited RH theorem therefore
makes these particular finite convolutions positive.
The source projections themselves are positive independently
of RH; positivity of their zero-dependent complex weights is
the additional issue. This is a concrete realization of \(B\)
and its kernel in the actual source Bessel model, not an
independent arithmetic derivation of those weights or a claim
that the historical period determinants are these finite sums.

\paragraph{Grothendieck framework for \fid{DET-12--15}.}
The signed Gram determinant and its positive determinant line are
finite Hermitian linear algebra; the positive frame determinants
are exterior-power determinants. Their precise algebraic home is
SGA~6, Expos\'e~V, Definition~2.1, Example~2.2(b), Lemma~2.2.1
and the trace logarithm (7.1.1) \cite{sga6}.
The passage from summable rank-one operators to the ordinary
Fredholm determinant belongs to cote~1, scans~78--80, with the
growth and genus discussion on scans~111--112 \cite{archive}.
The automorphic period and Bessel spaces carry the source's stated
analytic representation hypotheses. These arguments do not
identify them with the cohomology of a scheme over $\Spec\Z$,
an $\ell$-adic realization, or a crystal. The comparison with
SGA~5 trace formulas and Deligne's weights is made separately in
the finite-field section, where those geometric objects are
actually supplied. No EGA or FGA locator is asserted for an
arithmetic cohomology construction absent from this argument.
